% =====================================================================
%  Thème 3 — Relation fondamentale & formules d'addition — NIVEAU DIFFICILE
%  Corrigé
% =====================================================================
\documentclass[11pt,a4paper]{article}
\usepackage[utf8]{inputenc}
\usepackage[T1]{fontenc}
\usepackage{lmodern}
\usepackage[french]{babel}
\usepackage{amsmath,amssymb,mathtools}
\usepackage{icomma}
\usepackage{xcolor}
\usepackage{colortbl}
\usepackage{tikz}
\usepackage[most]{tcolorbox}
\usepackage{enumitem}
\usepackage{array}
\usepackage[a4paper,margin=2.2cm]{geometry}
\usetikzlibrary{arrows.meta,calc}

\definecolor{primary}{RGB}{214,51,108}
\definecolor{primarydark}{RGB}{140,20,64}
\definecolor{excol}{RGB}{0,135,90}

\DeclareMathOperator{\tg}{tg}
\DeclareMathOperator{\cotg}{cotg}
\newcommand{\degr}{^{\circ}}
\newcommand{\Z}{\mathbb{Z}}

\newcounter{exo}
\newcommand{\exo}[1]{\medskip\refstepcounter{exo}%
  \noindent{\large\bfseries\color{primarydark}Exercice \theexo}%
  \ \textcolor{primary}{\textbullet}\ \textbf{#1}\par\nopagebreak\smallskip}
\newcommand{\rep}{\textcolor{excol}{$\blacktriangleright$}\ }

\setlength{\parindent}{0pt}
\pagestyle{empty}

\begin{document}

\begin{center}
{\LARGE\bfseries\color{primary}Thème 3 — Relation fondamentale \& formules d'addition}\\[2pt]
{\color{primarydark}\textsc{Niveau Difficile — Corrigé détaillé}}
\end{center}
\hrule height 1pt
\medskip

\exo{La chaîne complète}
\begin{itemize}[leftmargin=1.4em,itemsep=2pt]
\item[\textbf{(a)}] \rep $\cos^2\alpha=1-\sin^2\alpha=1-\dfrac19=\dfrac89$. Au quadrant I $\cos>0$, donc
$\cos\alpha=\sqrt{\dfrac89}=\dfrac{2\sqrt2}{3}$.
\item[\textbf{(b)}] \rep $\sin2\alpha=2\sin\alpha\cos\alpha=2\cdot\dfrac13\cdot\dfrac{2\sqrt2}{3}
=\dfrac{4\sqrt2}{9}$ ; \ $\cos2\alpha=1-2\sin^2\alpha=1-\dfrac29=\dfrac79$.
\item[\textbf{(c)}] \rep $\sin\!\Big(2\alpha-\dfrac{\pi}{3}\Big)
=\sin2\alpha\cos\dfrac{\pi}{3}-\cos2\alpha\sin\dfrac{\pi}{3}
=\dfrac{4\sqrt2}{9}\cdot\dfrac12-\dfrac79\cdot\dfrac{\sqrt3}{2}=\dfrac{4\sqrt2-7\sqrt3}{18}$.
\item[\textbf{(d)}] \rep $\cos\!\Big(2\alpha+\dfrac{\pi}{3}\Big)
=\cos2\alpha\cos\dfrac{\pi}{3}-\sin2\alpha\sin\dfrac{\pi}{3}
=\dfrac79\cdot\dfrac12-\dfrac{4\sqrt2}{9}\cdot\dfrac{\sqrt3}{2}=\dfrac{7-4\sqrt6}{18}$.
\end{itemize}

\exo{Démontrer les formules de l'angle double}
\begin{itemize}[leftmargin=1.4em,itemsep=2pt]
\item[\textbf{(a)}] \rep Dans $\cos(a+b)=\cos a\cos b-\sin a\sin b$, on pose $b=a$ :
\[
\cos2a=\cos(a+a)=\cos a\cos a-\sin a\sin a=\cos^2a-\sin^2a.
\]
\item[\textbf{(b)}] \rep On remplace tour à tour à l'aide de $\cos^2a+\sin^2a=1$ :
\[
\cos2a=\cos^2a-\sin^2a=\cos^2a-(1-\cos^2a)=2\cos^2a-1,
\]
\[
\cos2a=\cos^2a-\sin^2a=(1-\sin^2a)-\sin^2a=1-2\sin^2a.
\]
\item[\textbf{(c)}] \rep En partant de $\cos2x=1-2\sin^2x$ (avec $a=x$), on isole $\sin^2x$ :
\[
2\sin^2x=1-\cos2x\quad\Longrightarrow\quad \sin^2x=\dfrac{1-\cos2x}{2}.
\]
\end{itemize}

\exo{Point sur le cercle et angle double}
\begin{itemize}[leftmargin=1.4em,itemsep=2pt]
\item[\textbf{(a)}] \rep $1+\tg^2\alpha=1+\dfrac94=\dfrac{13}{4}=\dfrac{1}{\cos^2\alpha}$, d'où
$\cos^2\alpha=\dfrac{4}{13}$. Quadrant II : $\cos<0$, donc
$\cos\alpha=-\dfrac{2}{\sqrt{13}}=-\dfrac{2\sqrt{13}}{13}$. Puis
$\sin\alpha=\tg\alpha\cdot\cos\alpha=\Big(-\tfrac32\Big)\Big(-\tfrac{2\sqrt{13}}{13}\Big)
=\dfrac{3\sqrt{13}}{13}$.
Ainsi $M=\left(-\dfrac{2\sqrt{13}}{13}\,;\,\dfrac{3\sqrt{13}}{13}\right)$.
\item[\textbf{(b)}] \rep
$\sin2\alpha=2\sin\alpha\cos\alpha=2\cdot\dfrac{3\sqrt{13}}{13}\cdot\Big(-\dfrac{2\sqrt{13}}{13}\Big)
=2\cdot\dfrac{-6\cdot13}{169}=-\dfrac{12}{13}$ ;
\[
\cos2\alpha=2\cos^2\alpha-1=2\cdot\dfrac{4}{13}-1=\dfrac{8}{13}-1=-\dfrac{5}{13}.
\]
\item[\textbf{(c)}] \rep Par le quotient :
$\tg2\alpha=\dfrac{\sin2\alpha}{\cos2\alpha}=\dfrac{-12/13}{-5/13}=\dfrac{12}{5}$.
Par la formule :
$\tg2\alpha=\dfrac{2\tg\alpha}{1-\tg^2\alpha}=\dfrac{2\cdot(-\frac32)}{1-\frac94}
=\dfrac{-3}{-\frac54}=\dfrac{12}{5}$. Les deux résultats coïncident.
\end{itemize}

\exo{Une formule de l'angle triple}
\begin{itemize}[leftmargin=1.4em,itemsep=2pt]
\item[\textbf{(a)}] \rep $\sin3a=\sin(2a+a)=\sin2a\cos a+\cos2a\sin a$.
\item[\textbf{(b)}] \rep On remplace $\sin2a=2\sin a\cos a$ et $\cos2a=1-2\sin^2a$ :
\[
\sin3a=(2\sin a\cos a)\cos a+(1-2\sin^2a)\sin a=2\sin a\cos^2a+\sin a-2\sin^3a.
\]
Avec $\cos^2a=1-\sin^2a$ :
\[
\sin3a=2\sin a(1-\sin^2a)+\sin a-2\sin^3a=2\sin a-2\sin^3a+\sin a-2\sin^3a=3\sin a-4\sin^3a.
\]
\item[\textbf{(c)}] \rep Avec $\sin a=\dfrac12$ :
$\sin3a=3\cdot\dfrac12-4\cdot\dfrac18=\dfrac32-\dfrac12=1$, ce qui est bien $\sin90\degr=1$.
\end{itemize}

\end{document}
